\documentclass[11pt,a4paper]{article}
\usepackage[utf8]{inputenc}
\usepackage[T1]{fontenc}
\usepackage{amsmath,amssymb,amsthm}
\usepackage[margin=1in]{geometry}
\usepackage{booktabs}
\usepackage{longtable}
\usepackage{url}
\usepackage[hidelinks]{hyperref}
\usepackage{enumitem}

\theoremstyle{plain}
\newtheorem{proposition}{Proposition}
\newtheorem{theorem}[proposition]{Theorem}
\newtheorem{lemma}[proposition]{Lemma}
\theoremstyle{definition}
\newtheorem{hypo}{Hypothesis}
\renewcommand{\thehypo}{$\beta$}

\newcommand{\Fc}{F_{c}}
\newcommand{\al}{\alpha}

\title{Exact Realizer Floors for One- and Two-Defect\\
Valuation Words in the Accelerated $3x+1$ Map}
\author{Elias de Jesus}
\date{22 October 2026}

\begin{document}
\maketitle

\begin{abstract}
For the accelerated Syracuse map $T(x)=(3x+1)/2^{v_2(3x+1)}$ on odd integers we
give the exact least realizer of a two-defect valuation word
$w=1^{u}q_a1^{v}q_b1^{t}$ in closed form, and an unconditional
per-configuration floor $\Fc(w)$ obtained by combining truncation with
orbit-advance at one-defect granularity. We record the modulus convention
precisely: the exact-valuation realizer set is one residue class modulo
$2^{S+1}$, not $2^{S}$. We prove that the anchor denominator is $3^{n}$ on the
interior stratum $u,v,t\ge 1$ and identify the boundary stratum $t=0$ where it
drops, the extreme witness being the $x=1$ cycle word. We then sharpen the floor to $\Fc'$ using the fact that the second defect closes the
first episode, and exhibit an infinite family on which that sharpening gains an
unbounded number of bits. We then reduce the \emph{uniform} two-defect floor to
a single named one-defect import $(\beta)$, with the threshold $\beta>1/2$
exact at every scale, and record two negative results with computed witnesses.

\medskip\noindent\textbf{Non-claims, stated up front.}
Nothing here bears on the Collatz conjecture and nothing excludes divergent
orbits; every statement concerns least realizers of finite valuation words.
There is no formal verification: ``verified'' means exhaustive or exact finite
computation on a stated grid. Theorem~\ref{thm:fc} is a per-configuration
certificate and is \emph{not} a uniform family floor. The hypothesis $(\beta)$
is open. A limited prior-art audit was carried out (\S\ref{sec:lit}) but is not
exhaustive, so novelty is not established; branches (a) and (b) of
Theorem~\ref{thm:fc} are standard cylinder facts and
Proposition~\ref{prop:class} restates a correction published earlier by the
author. What is claimed here is the closed-form per-configuration evaluation,
the closed-episode branch~(c) with its unbounded gain
(Proposition~\ref{prop:unbounded}), and the $\beta$-reduction.
\end{abstract}

\section{Epistemic tiers}

\begin{center}
\begin{tabular}{@{}llp{6.1cm}@{}}
\toprule
Tier & Meaning & Used for \\
\midrule
T1 & Exact integer arithmetic; proof & Props.~1--2, Thms.~3,~5,
Lem.~4, all grid counts \\
T2 & Exact integers, \texttt{math.log2} on exact inputs & $\delta$ values
(\S\ref{sec:neg}); Pearson estimates (\S\ref{sec:limits}) \\
T3 & Interpretation & framing remarks, \S\ref{sec:limits} items 4--6 \\
\bottomrule
\end{tabular}
\end{center}

\noindent\texttt{mpmath} was unavailable in the working environment, so the T2
tier is served by exact integers rather than arbitrary-precision floats, which
is stronger than required. No floating-point value decides any assertion.

\section{The realizer class and the modulus convention}\label{sec:class}

Throughout $\al=\log_2 3=1.5849625\dots$.
A valuation word is $w=(v_1,\dots,v_n)$ with $v_k\ge 1$; set
$S_k=v_1+\cdots+v_k$, $S=S_n$, $|w|=n$. The affine identity is
$2^{S_k}x_k=3^kx+C_k$ with append recurrence $C'=3C+2^{S}$, equivalently
$C_w=\sum_{k=1}^{n}3^{\,n-k}2^{S_{k-1}}$.

\begin{proposition}\label{prop:class}
An odd $m$ realizes $w$ with exact valuations if and only if
$3^nm+C_w\equiv 2^{S}\pmod{2^{S+1}}$. The realizer set is exactly one residue
class modulo $2^{S+1}$; write $m(w)$ for its least positive member.
\end{proposition}

\begin{proof}
Two steps. (a) The top divisibility $2^{S_n}\mid 3^nm+C_n$ implies every
earlier one: from $C_n=3C_{n-1}+2^{S_{n-1}}$,
$3^nm+C_n=3(3^{n-1}m+C_{n-1})+2^{S_{n-1}}$, and $2^{S_{n-1}}$ divides the left
side and the trailing term, hence $3(3^{n-1}m+C_{n-1})$ and so
$3^{n-1}m+C_{n-1}$. The same computation gives $2^{v_n}\mid 3x_{n-1}+1$, so
each $x_k$ with $k<n$ is automatically odd. (b) The one remaining condition is
that $x_n$ be odd, worth exactly one further bit; it upgrades the integrality
modulus $2^{S}$ to $2^{S+1}$.
\end{proof}

The \emph{integrality} class alone is modulo $2^{S}$; the exact-valuation class
is modulo $2^{S+1}$. The distinction is the difference between a realizer and
the least realizer, and it is the source of the reconciliation in
\S\ref{sec:tworun}.

\paragraph{Verification.}
An independent implementation sharing no code with the first builds the class by
backward recursion, $x\equiv 3^{-1}(2^{v_1}a'-1)\pmod{2^{v_1+M'}}$, from the
base ``the empty word is realized by every odd $x$''. The modulus $2^{S+1}$
emerges with no input, and a full scan of \emph{every} odd $m<2^{S+1}$ on $243$
words confirms a singleton class with $M=S+1$ in every case: $0$ failures.

\section{One defect, and a notation reconciliation}\label{sec:tworun}

The conjugation $y=x+1$ linearises a run of ones: $v_k=1$ gives
$x_k+1=3(x_{k-1}+1)/2$, so $x_k+1=(3/2)^k(x_0+1)$. A leading run of $u$ ones
forces $2^{u}\mid m+1$, and a following letter $q\ge 2$ forces
$2^{u+1}\,\|\,m+1$.

\begin{proposition}\label{prop:tworun}
Let $m+1=2^{u+1}A$ with $A$ odd, $n=q+v$, and let
$y=(3^{u+1}A-1)/2^{\,q-1}$ be the odd orbit value reached after $1^{u}q$.
Then the run of $v$ ones is governed by $v_2(y+1)$, and the two natural
conditions are \emph{distinct}:
\begin{center}\small
\begin{tabular}{@{}lll@{}}
\toprule
condition on the orbit & meaning & congruence\\
\midrule
$v_2(y+1)\ge v+1$ & realizes the finite word $1^{u}q1^{v}$
  & $3^{u+1}A\equiv 1-2^{q-1}\pmod{2^{n}}$\\
$v_2(y+1)=v+1$ & \dots\ and the next valuation is $\ge2$
  & $3^{u+1}A\equiv 1-2^{q-1}+2^{n}\pmod{2^{n+1}}$\\
\bottomrule
\end{tabular}
\end{center}
Write $m(P)$ for the least realizer under the first condition and
$m_{\mathrm{closed}}(P)$ for the least under the second.
\end{proposition}

\begin{proof}
After $1^{u}$ the orbit is at $2\cdot3^{u}A-1$, so
$3x+1=2(3^{u+1}A-1)$ and $v_2=q$ forces $v_2(3^{u+1}A-1)=q-1$, giving
$x_{u+1}=y$. A run of $v$ ones needs $2^{v}\mid y+1$ (conjugate $y\mapsto y+1$,
under which the run is $y\mapsto 3y/2$), and the final value
$3^{v}(y+1)/2^{v}-1$ is odd exactly when $2^{v+1}\mid y+1$. For the next
valuation, $3x_n+1=2\bigl(3^{v+1}(y+1)/2^{\,v+1}-1\bigr)$, which is divisible
by $4$ precisely when $2^{v+1}\,\|\,y+1$.
\end{proof}

\paragraph{The source's finer modulus is not an error.}
The imported note states the mod-$2^{n+1}$ form. That is the \emph{closed
episode} condition, a different and equally legitimate specification---the end
of a whole episode rather than an arbitrary finite prefix---not a typo for
$2^{n}$. Both are used below: Proposition~\ref{prop:tworun}'s first row for the
general theory, its second row in Theorem~\ref{thm:fc}\,(c).

\paragraph{The two realizers can differ by many bits.}
$m_{\mathrm{closed}}(P)\in\{m(P),\,m(P)+2^{S_P+1}\}$, so the two differ by at
most one \emph{modulus}---but their \emph{heights} can differ by far more than
one bit, since $m(P)$ may sit well below $2^{S_P+1}$. At $u=11$, $q=5$, $v=4$:
$m(P)=4095$ while $m_{\mathrm{closed}}(P)=2\,101\,247=4095+2^{21}$, so
$\log_2\bigl(m_{\mathrm{closed}}/m\bigr)=9.0032$. The measured supremum grows
with the grid ($7.4237$ for $u,v\le10$; $9.0032$ for $u,v\le14$; $9.4294$ for
$u,v\le18$), so this is not an $O(1)$ correction. One extra congruence bit is
not one extra bit of height, and any estimate comparing the two must use the
height, not the modulus count.

\paragraph{Placement reduction.}
Minimality of $A$ is governed by where $c\,3^{-(u+1)}$ falls in
$\mathbb{Z}/2^{n}$; dually, by the ternary digits of $2^{-n}$ modulo
$3^{u+1}$. This is the one-dimensional case of the obstruction in
\S\ref{sec:limits}.

\section{Two defects: closed form and denominator}\label{sec:r1}

Let $w=1^{u}q_a1^{v}q_b1^{t}$, $n=u+v+t+2$, $S=u+v+t+q_a+q_b$.

\begin{theorem}\label{thm:r1}
$C_w=3^{n}-3^{\,v+t+1}2^{\,u+1}+3^{\,v+t+1}2^{\,u+q_a}
-3^{\,t}2^{\,u+q_a+v+1}+3^{\,t}2^{\,u+q_a+v+q_b}-2^{S}.$
\end{theorem}

Equivalently, by Proposition~\ref{prop:class},
\[
m(w)\equiv -1-3^{-(u+1)}2^{u+1}\!\left(2^{q_a-1}-1\right)
-3^{-(u+v+2)}2^{\,u+q_a+v+1}\!\left(2^{q_b-1}-1\right)\pmod{2^{S+1}},
\]
which exhibits $2^{u+1}\,\|\,m+1$, both cofactors being odd. At $q_a=q_b=1$ the
formula collapses to $3^{n}-2^{n}$, the all-ones value. Verified: $0$
mismatches out of $19\,773$ against the append recurrence.

\begin{lemma}\label{lem:denom}
The anchor is $3^{-n}(2^{S}-C_w)$, with reduced denominator
$D=3^{\,\max(0,\;n-v_3(2^{S}-C_w))}$. On the interior stratum $u,v,t\ge1$,
$v_3(2^{S}-C_w)=0$ and hence $D=3^{n}$.
\end{lemma}

\begin{proof}
With $t\ge 1$ every monomial of $C_w$ except $-2^{S}$ carries a factor $3$, so
$C_w\equiv -2^{S}\pmod 3$ and $2^{S}-C_w\equiv 2^{S+1}\not\equiv 0\pmod 3$.
\end{proof}

\paragraph{The denominator is representative-dependent.}
Both one-sided lifts pin one representative each; the canonical Lemma-H datum
is the \emph{balanced} lift $N^{*}\equiv-C\pmod{2^{S}}$ with
$|N^{*}|\le 2^{S-1}$, giving $D=3^{\,n-v_3(N^{*})}$. On the family $u,v,t\in1..12$ with $q_a,q_b\in\{2,3,4\}$ ($15\,552$
configurations), verified independently: both
one-sided representatives are $3$-coprime everywhere ($0/15\,552$ each); the
balanced lift reduces on $5263/15\,552=0.3384$, with $v_3(N^{*})$ following a
geometric$(2/3)$ law to within sampling error:

\begin{center}\small
\begin{tabular}{@{}lrrrrrrrrrr@{}}
\toprule
$k$ & 0 & 1 & 2 & 3 & 4 & 5 & 6 & 7 & 8 & 9\\
\midrule
observed & 10289 & 3478 & 1199 & 417 & 123 & 32 & 9 & 3 & 1 & 1\\
geom$(2/3)$ & 10368 & 3456 & 1152 & 384 & 128 & 42.7 & 14.2 & 4.7 & 1.6 & 0.5\\
\bottomrule
\end{tabular}
\end{center}

\noindent The dangerous-window probe gives median $\log_2m/(S+1)$ of $0.9615$
at $v_3(N^{*})=0$ versus $0.9620$ at $v_3(N^{*})>0$, a difference of $0.00046$. On these grids the reduction is therefore
\textbf{consistent with the geometric benchmark and shows no predictive signal
about $m$}. This is a finite measurement, not a proof that no relationship
exists.

The premise $\gcd(C_w,3^n)=1$ holds on all $19\,773$ configurations but does
\emph{not} imply $D=3^n$, since $D$ is governed by $\gcd(2^{S}-C_w,3^{n})$.
Stratifying over $|w|\le 44$, reductions $v_3>0$ occur on $1176$
configurations with $t=0$ and on $0$ with $t\ge1$: the reduction is confined to
$t=0$ exactly. The truncation at $0$ is necessary, not cosmetic: $v_3$ may exceed $n$. For
$w=(6,2)$, $n=2$, $S=8$, $C_w=67$ and $2^{S}-C_w=189=3^{3}\cdot7$, so
$v_3=3>n$; the anchor is $189/9=21$ with reduced denominator $1$, not
$3^{-1}$. (On a grid with $u,v,t\le6$ and $q_a,q_b\in\{2,\dots,7\}$ this occurs on
$4/12\,348$ configurations, all with $v=t=0$; it is invisible on a grid
restricted to $q_a,q_b\in\{2,3,4\}$.) The extreme witness is the $x=1$ cycle word
$w=(2,2)$, where
$m(w)=1$, $2^{S}-C_w=9$, $v_3=2$ and $D=1$: reduced denominators mark cycle
proximity, consistent with the cycle-denominator identification of Lagarias~\cite{lagarias1990}. Accordingly the family is stated as $u,v,t\ge1$ and the $t=0$ and
$v=0$ strata are carved out explicitly.

\section{The suffix-combination floor}

Let $P=1^{u}q_a1^{v}$, $P'=1^{u}q_a$, $Q=1^{v}q_b1^{t}$.

\begin{theorem}\label{thm:fc}
(a) $m(w)\ge m(P)$;\quad
(b) $m(w)\ge\left\lceil\left(2^{S_{P'}}m(Q)-C_{P'}\right)/3^{\,u+1}\right\rceil$;\quad
(c) $m(w)\ge m_{\mathrm{closed}}(P)\ \ge m(P)$.
Hence $\Fc'(w):=\max\{(b),(c)\}$ is an unconditional floor, computable in
closed form for each configuration, and $\Fc'\ge\Fc:=\max\{(a),(b)\}$.
\end{theorem}

\begin{proof}
\emph{(a) Truncation:} any realizer of $w$ realizes the prefix $P$, so $m(w)$
lies in the class of $P$. \emph{(b) Orbit-advance:} $x_{u+1}$ realizes $Q$, so
$x_{u+1}\ge m(Q)$, and $m=(2^{S_{P'}}x_{u+1}-C_{P'})/3^{u+1}$ is increasing in
$x_{u+1}$. \emph{(c)} In $w=1^{u}q_a1^{v}q_b1^{t}$ the prefix $P$ is followed
by $q_b\ge2$, so every realizer of $w$ realizes $P$ \emph{as a closed
episode}; apply the second row of Proposition~\ref{prop:tworun}.
\end{proof}

\paragraph{The closed branch is a strict improvement.}
Branch (c) is available only because the second defect closes the first
episode---it is not an alternative convention but extra information the
two-defect structure supplies for free. Measured on $u,v,t\in1..11$
($15\,552$ configurations): $\Fc'$ is \emph{never} worse than $\Fc$ and is
strictly better on $5442$ ($35.0\%$), with mean gain $1.649$ bits, median
$1.414$, maximum $4.790$. Truncation binds in $64.2\%$ of configurations. On that grid the
improvement is small---the balanced median rises from $0.6826$ to $0.6903$ per
unit $|w|$ against an empirical $1.0341$---but \textbf{it is not $O(1)$ in
general}, and an earlier draft's claim that it was has been withdrawn.

\begin{proposition}\label{prop:unbounded}
The gain $\log_2(\Fc'/\Fc)$ is unbounded on the two-defect family. For each
$v\ge1$ choose an even $e$ with $3^{e}\equiv-15\pmod{2^{\,v+6}}$ and $e$ large
enough that branch (b) is dominated. For $w=1^{\,e-1}5\,1^{v}2\,1$,
\[
\Fc(w)=2^{e}-1,\qquad \Fc'(w)=2^{e}-1+2^{\,e+v+5},\qquad
\log_2\frac{\Fc'(w)}{\Fc(w)}>v+5 .
\]
\end{proposition}

\begin{proof}
\emph{Existence of arbitrarily large even $e$.} For $k\ge3$ the group
$(\mathbb{Z}/2^{k})^{*}$ is $\{\pm1\}\times\langle3\rangle$ with
$\operatorname{ord}(3)=2^{\,k-2}$, so $\langle9\rangle=\langle3^{2}\rangle$ has
order $2^{\,k-3}$ and is contained in $\{x\equiv1\bmod 8\}$, a set of the same
cardinality; hence $\langle9\rangle=\{x\equiv1\bmod 8\}$ exactly. Since
$-15\equiv1\pmod 8$, the congruence $3^{e}\equiv-15\pmod{2^{\,v+6}}$ is
solvable, and because $3^{e}\equiv1\pmod 8$ iff $e$ is even, every solution is
even. The solutions form the progression $e\equiv e_0 \pmod{2^{\,v+4}}$, so
arbitrarily large even solutions exist.

\emph{The two floors.} Put $u=e-1$, $q_a=5$, so $P=1^{\,e-1}5\,1^{v}$,
$n=q_a+v=v+5$ and $S_P=e+v+4$. Taking $A=1$ in
Proposition~\ref{prop:tworun} gives $m=2^{e}-1$, and the first row requires
$3^{e}\equiv1-2^{4}=-15\pmod{2^{\,v+5}}$, which our congruence implies; as
$A=1$ is the least positive odd value, $m(P)=2^{e}-1$. The second row requires
$3^{e}\equiv-15+2^{\,v+5}\pmod{2^{\,v+6}}$, which our congruence
\emph{contradicts}; so $A=1$ is not closed, and
$m_{\mathrm{closed}}(P)=m(P)+2^{S_P+1}=2^{e}-1+2^{\,e+v+5}$. This is exactly
why the hypothesis is imposed modulo $2^{\,v+6}$ and not $2^{\,v+5}$: the finer
modulus simultaneously selects $A=1$ for the open word and rejects it for the
closed one.

\emph{Branch (b) is dominated.} With $t=1$, $Q=1^{v}2\,1$ and $m(Q)$ depend only
on $v$, while $S_{P'}=u+q_a=e+4$. Hence branch (b) is at most
$2^{\,e+4}m(Q)/3^{\,e}$, and its ratio to $2^{e}-1$ is
$O\bigl(16\,m(Q)\,3^{-e}\bigr)\to0$. So for all large $e$ in the progression the
truncation branch binds in both $\Fc$ and $\Fc'$, giving
$\log_2(\Fc'/\Fc)=\log_2\bigl(1+2^{\,v+5}\cdot\tfrac{2^{e}}{2^{e}-1}\bigr)>v+5$.
\end{proof}

Each step was checked numerically: $\langle9\rangle=\{x\equiv1\bmod8\}$ for
$2^{4}\dots2^{12}$; the solution progressions and their parity for
$v=1,\dots,5$; the open/closed selection of $A=1$; and the domination of
branch~(b) along the progression. Verified for $v=1,\dots,7$: both closed forms match the computed floors
exactly, and the gain is $6.0227, 7.0116, 8.0060, 9.0032, 10.0014, 11.0007,
12.0004$, decreasing to $v+5$ from above. The smallest instance is already
instructive: $w=1^{11}5\,1^{4}2\,1$ has $\Fc=4095$, $\Fc'=2\,101\,247$ and
$m(w)=14\,684\,159$, a gain of $9.0032$ bits with both floors valid. Branch (c)
therefore sharpens the per-configuration certificate by an unbounded amount on
suitable words, while leaving \S\ref{sec:beta} and every asymptotic claim
untouched (those concern the balanced regime, where the gain is small).

\paragraph{Framing (T3).}
$\Fc$ is \textbf{not a new mechanism}: it is orbit-advance plus truncation,
evaluated at one-defect granularity---truncating to the full one-defect prefix
$P$ and using its exact $m(P)$ rather than to a maximal all-ones run. The
novelty claimed is only closed-form per-configuration evaluability and the
$\beta$-reduction of \S\ref{sec:beta}.

\paragraph{Independence certificate.}
$\Fc\le m(w)$ with $0$ violations out of $19\,773$, and $\Fc'\le m(w)$ with $0$ out of $15\,552$. Against the
all-ones-run floor: better on $19\,252$, tied on $521$, worse on $0$;
dominant, strictly better on $97.4\%$. The second, independent implementation
of the same specification reports $450/451$ with one exact tie at a degenerate
$|w|=32$, $v=1$ configuration---the same qualitative result on a different
grid. That package was run in this environment and reproduced its balanced
medians $0.337/0.671/1.015$ exactly, against $0.683/0.336/1.034$ here; the
differences are grid composition.

\section{The uniform reduction}\label{sec:beta}

Theorem~\ref{thm:fc} is per-configuration. A uniform family floor needs a
uniform one-defect input.

\begin{hypo}\label{hyp:beta}
$\log_2 m(1^{U}q1^{V})\ \ge\ \beta\cdot S(1^{U}q1^{V})$ for all
$U\ge1$, $V\ge1$, $q\ge2$.
\end{hypo}

\paragraph{The domain restriction is necessary.}
With $V=0$ allowed the hypothesis is false for any $\beta>1/2$: the word
$(1,4)$, i.e.\ $U=1$, $q=4$, $V=0$, has $m=3$ and $S=5$, giving
$\log_2m/S=0.3170$. Allowing $U=V=0$ is worse still---$w=(2)$ has $m=1$ and
ratio $0$. On the interior $U,V\ge1$, with $q\in\{2,\dots,7\}$, the measured minimum
is $0.5714$, attained
at $(U,q,V)=(11,5,5)$, and it is \emph{stable}: the same minimiser and value
on grids $U,V\le20$, $\le30$ and $\le45$ ($0.5977$ for $U,V\le10$). So the
interior family is consistent with $\beta>1/2$, with a margin of about $0.07$
and no observed drift toward $1/2$---evidence for the hypothesis, not a proof
of it. The two-defect application only ever invokes
Hypothesis~\ref{hyp:beta} on the interior, since $P=1^{u}q_a1^{v}$ with
$u,v\ge1$.

Normalisation, stated explicitly: balanced $u=v=t=k$, $q_a=q_b=2$, so
$S_P=2k+2$, and the comparator is the all-ones-run floor of $k+1$ bits.
\begin{itemize}[nosep]
\item \emph{Truncation branch:} $\beta(2k+2)>k+1\iff\beta>(k+1)/(2k+2)=1/2$,
exactly $1/2$ for \emph{every} $k$, not merely asymptotically.
\item \emph{Advance branch:} $\beta>(\al(k+1)-1)/(2k+2)\to\al/2=0.792481$,
rising $0.542\to0.792$ across $k=1,\dots,5000$.
\end{itemize}
Empirically $\beta$ for $1^{k}21^{k}$ is $0.9167$ $(k=5)$, $0.9286$ $(20)$,
$0.9672$ $(60)$, $1.0033$ $(150)$. \textbf{Hypothesis~\ref{hyp:beta} with
$\beta>1/2$ is the single question to which this note reduces the uniform
two-defect problem. It is open, and it is an effective question.}

\section{Negative results}\label{sec:neg}

\paragraph{Merge bridge: failed.}
Canonical image $m_2(u,q_a+q_b,v+t)$. At $(u,q_a,v,q_b,t)=(22,4,2,2,6)$,
$|w|=32$:

\begin{center}
\begin{tabular}{@{}lccc@{}}
\toprule
image & bits & $\delta$ (bit length) & $\delta$ (true $\log_2$, T2)\\
\midrule
canonical $m_2(u,q_a+q_b,v+t)$ & 37 & $+0.3125$ & $\mathbf{+0.3209}$\\
$m_2(u,q_a,v+t)$ & 34 & $+0.2188$ & $+0.2248$\\
$m_2(u,q_a+q_b-1,v+t)$ & 35 & $+0.2500$ & $+0.2508$\\
\bottomrule
\end{tabular}
\end{center}

\noindent The canonical variant reproduces the second implementation's
$\delta=0.321$ to four decimals, provided $\delta$ is measured by true $\log_2$
rather than bit length; that cross-engine agreement is the engines-agree
certificate. Suprema of $\delta$ per unit $|w|$: $+0.8662$ at $(2,2,1,4,1)$,
$|w|=6$ on the small full grid and on its interior; $+0.2093$ at
$(21,2,21,3,0)$, $|w|=44$ for $|w|\ge32$. The small-grid suprema are attained
at the smallest admissible words and are boundary artefacts, so the
$|w|\ge 32$ figure is the one bearing on asymptotics.

\paragraph{Deletion bridge: dominated} by Theorem~\ref{thm:fc}---it requires
$\beta\gtrsim0.56$, strictly worse than $1/2$, and additionally needs a
uniform-cost hypothesis that $\Fc$ does not.

\paragraph{Why no bridge beats $\Fc$ (T3).}
A bridge posits $m(v)\le A\,m(w)^{q}2^{\delta|w|}$, which rearranges to
$m(w)\ge\bigl(m(v)/(A2^{\delta|w|})\bigr)^{1/q}$---\emph{the bridge's own
conclusion}. So that condition is never weaker than the floor it produces, and
a bridge is non-circular only when it is established by a mechanism indifferent
to the size of $m(w)$. Exactly one such mechanism exists: a realizer realizes
every contiguous factor of $w$, directly or after orbit advance. Those are the
two branches of $\Fc$. Merge and deletion images are not factors of $w$.

\section{Calibration}

Medians of bits per $|w|$, balanced regime, $|w|\ge62$, $369$ configurations:
all-ones-run floor $0.336$; $\Fc$ $0.683$; $\Fc'$ $\mathbf{0.690}$;
empirical $1.034$.
Per regime at $|w|\approx122$ (bits):

\begin{center}
\begin{tabular}{@{}lrrrrrl@{}}
\toprule
regime & trunc & adv & $\Fc$ & run floor & empirical & binds\\
\midrule
balanced & 82 & 59 & \textbf{82} & 41 & 124 & truncation\\
leading-heavy & 122 & 1 & \textbf{122} & 119 & 122 & truncation\\
middle-heavy & 119 & 122 & \textbf{122} & 119 & 125 & advance\\
trailing-heavy & 5 & 119 & \textbf{119} & 119 & 124 & advance\\
\bottomrule
\end{tabular}
\end{center}

\noindent $\Fc$ is near-sharp in the lopsided regimes ($0.93$--$1.00$ of the
empirical exponent) and loses about a third of it when balanced. Under
Hypothesis~\ref{hyp:beta} with $\beta\approx0.95$ the conditional balanced rate
is $\approx0.633$ per unit word length, so $(\beta)$ does \emph{not} by itself
close the balanced gap; it closes the gap to the run-floor comparator.


\section{Relation to prior work}\label{sec:lit}

A literature check was performed for this revision. It is not exhaustive, and
absence from it is not evidence of novelty.

\paragraph{The affine identity is classical.}
$2^{S_k}x_k=3^{k}x+C_k$ with $C_w=\sum_k 3^{\,n-k}2^{S_{k-1}}$ is
Proposition~4 of B\"ohm and Sontacchi~\cite{bohmsontacchi} (1978), in the
two-branch form $u^{n}(x)=\bigl(3^{m}x+\sum_{k<m}3^{\,m-k-1}2^{v_k}\bigr)/2^{n}$;
their Proposition~5 is the cycle closed form. \emph{Verified first-hand.}

\paragraph{Proposition~\ref{prop:class} is the author's own prior correction.}
That the exact realizers of a valuation word form a single residue class modulo
$2^{S+1}$ rather than the Terras--Everett modulus $2^{S}$ was established in
\cite{dejesus-erratum} (11 August 2026), explicitly as an erratum to the
Terras--Everett residue class, and is carried as Proposition~5.2 of
\cite{dejesus-eoc}. Proposition~\ref{prop:class} restates it here, with a proof,
only because the $2^{S}$ versus $2^{S+1}$ distinction drives
\S\ref{sec:tworun}; no new priority is claimed for it in this note.
The same statement appears independently in
Akhtar~\cite{akhtar2026} (19 August 2026) as Theorem~9, with the explicit class
\[
r(k)\;\equiv\;-\sum_{j=0}^{m}2^{K_j}3^{-(j+1)}\pmod{2^{K+1}},
\qquad K_j=\textstyle\sum_{i<j}k_i ,
\]
derived from the parity-vector bijection. \emph{Full text consulted.} That
expression and the form used here, $3^{-n}(2^{S}-C_w)\bmod 2^{S+1}$, are
equivalent: they differ by $2^{K+2}\cdot(\text{odd})\equiv0$, verified on $1944$
two-defect words with $0$ disagreements.

Akhtar's concerns are otherwise disjoint from \S5--\S6. He uses the least
nonnegative representative only in normalized form, $r/M$ with $M=2^{K+1}$, as a
placement diagnostic (his Table~4); he does not bound least realizers from
below, give closed forms for them on structured word families, or use
prefix-truncation and orbit-advance as floors. His ``prefix'' quantities are
prefix sums $K_i$ and worst-prefix excursions, not truncations of the word to a
shorter one. The closed-episode condition of
Proposition~\ref{prop:tworun} does not appear.

\paragraph{Parity vectors.}
For the shortcut map, the bijection between parity vectors of length $k$ and
residue classes modulo $2^{k}$ is due to Terras~\cite{terras1976} (1976) and
independently Everett~\cite{everett1977} (1977); Lagarias's survey
\cite{lagarias2010} credits both and points to \cite[Theorem~B]{lagarias1985}
for the precise statement, in the form of the finitely additive invariant
measure assigning mass $2^{-n}$ to each arithmetic progression modulo $2^{n}$.
\emph{Verified through the survey, whose full text was consulted; the 1976,
1977 and 1985 originals were not.}

Rozier and Terracol~\cite{rozierterracol} treat parity vectors for the map
$n\mapsto(3n+1)/2$ on odd $n$, with $V_j(n)\in\{0,1\}^{j}$, introducing a
partial order on parity vectors and studying paradoxical sequences.
\emph{Full text consulted (v5, 17 May 2026).} That is a different object from
the present setting: their parity vectors are binary words of length $j$ with
modulus $2^{j}$, whereas a valuation word here records the full exponents
$v_2(3x+1)$ and carries modulus $2^{S+1}$. They do not treat least elements of
such classes.

\paragraph{The leading-run spine.}
That a leading run of $u$ ones forces $2^{u+1}\mid m+1$, giving
$m(1^{u})=2^{u+1}-1$, is the ``modular resistance spine'' of
Kayadibi~\cite{kayadibi2026}, who shows $n=2^{m}q-1$ implies
$T^{j}(n)=3^{j}2^{\,m-j}q-1$ for $0\le j\le m-1$ and hence a deterministic
lower bound on first-descent time (Definition~3.1, Lemma~3.3,
Propositions~3.4 and~3.8). \emph{Full text consulted.} That paper does not
treat least realizers or the modulus $2^{S+1}$.

\paragraph{Extremal $C_w$ bounds.}
Sharp bounds on $C_w$---relevant to the anchor sizes here, though not to
anything we prove---are B\"ohm--Sontacchi Proposition~6, rediscovered as
Lemma~4 of Halbeisen and Hungerb\"uhler~\cite{hh1997}. \emph{Verified
first-hand.}

\paragraph{What the check did not find.}
No prior treatment was found of: the closed-episode/finite-word distinction of
Proposition~\ref{prop:tworun} and its consequence Theorem~\ref{thm:fc}\,(c);
the exact two-defect closed form (Theorem~\ref{thm:r1}); or the
$\beta$-reduction of \S\ref{sec:beta}. In particular Lagarias's survey
\cite{lagarias2010}, which covers the parity-vector framework and its
literature, contains no treatment of \emph{least} elements of parity-vector
classes, and neither does \cite{akhtar2026}, the one work found to state
Proposition~\ref{prop:class} independently. This is a negative search result,
not a novelty claim, and the audit remains non-exhaustive; but every work cited
here has now been read in full or, for three classical items, through
Lagarias's survey.

\section{Limitations}\label{sec:limits}

\begin{enumerate}[nosep,leftmargin=*]
\item Nothing here bears on the Collatz conjecture; nothing excludes divergent
orbits.
\item No formal verification of any kind.
\item Per-configuration $\ne$ uniform; the uniform statement needs
Hypothesis~\ref{hyp:beta}.
\item The residual above $\Fc$ is an incomplete-window statement: it requires
that $-C_w3^{-n}\bmod 2^{S+1}$ not be anomalously small. For two defects this
is a \emph{two-dimensional} pointwise window problem for
$\langle3\rangle\le(\mathbb{Z}/2^{M})^{*}$---simultaneous dyadic windows for
$3^{-a}$ and $3^{-a-b}$---whose one-dimensional case is \S\ref{sec:tworun}.
Character and complete-sum methods are provably unavailable there, $2$ being a
primitive root modulo $3^{b}$, and the two-dimensional case is strictly
stronger.
\item Effective methods degrade categorically with a second defect: the form
has three logarithms with heights $(3,A,Z)$,
$Z=3^{v+1}(2^{q_a-1}-1)+2^{q_a+v}(2^{q_b-1}-1)$, where $\log_2Z=\Theta(v)$
exactly and $Z$ is never $\{2,3\}$-smooth ($19\,773/19\,773$). The recoverable
exponent fraction therefore acquires a factor $1/v$ and tends to $0$ for every
fixed constant, except for bounded middle runs. The three-logarithm constant is
\textbf{not verified at source}, so no crossover is asserted.
\item The residual behaves as an independent residual. The published
statistic---$(\log_2m-\Fc)/|w|$ against P-anchor fullness
$\log_2m(P)/(u+q_a+v+1)$ at fixed $|w|=122$---reproduces exactly here,
$r=-0.179$, $n=100$. That design is underpowered: over $40$ replicates $r$
ranges over $[-0.372,+0.042]$, sd $0.111$, with $|r|>0.197$ firing in $21/40$,
so neither magnitude nor sign is stable. At $n=20\,000$ the same statistic
gives $r=-0.1673$, nonzero against the $0.0139$ threshold and close to the
published estimate. The correlation is real but explains $r^{2}=2.8\%$ of
residual variance, so the conclusion holds on grounds of \emph{small effect
size}, not statistical insignificance.
\item A literature check was performed (\S\ref{sec:lit}) but is not
exhaustive; absence from it is not evidence of novelty.
Proposition~\ref{prop:class} is the author's own prior result
(\cite{dejesus-erratum}) and is not claimed as new here.
\end{enumerate}

\section{Reproduction}

All code and data are public at
\url{https://github.com/innerlightr-wq/realizer-floors}.
Python 3.14.7, standard library only; no \texttt{numpy}, \texttt{sympy} or
\texttt{mpmath}.

\begin{verbatim}
python3 src/verify_note_numbers.py > data/note_numbers.txt
python3 src_indep/indep_realizer.py          # independent class, R1, denom
python3 src_indep/test_R2_strata.py          # denominator by stratum
python3 src_indep/test_R2_balanced_lift.py   # balanced lift + geom(2/3)
python3 src_indep/test_R3.py                 # F_c validity and dominance
python3 src_indep/test_R4_R5.py              # beta thresholds, merge witnesses
python3 src_indep/test_R5_canonical.py       # canonical merge image
python3 src_indep/test_R6_their_stat.py      # R6 power, published stat
python3 src/two_run_crosscheck.py            # two-run import cross-check
\end{verbatim}

Every number printed above is regenerated by \texttt{verify\_note\_numbers.py}.
Superseded scripts are retained but not listed: \texttt{test\_R6\_power.py}
(a correlation on a different statistic, withdrawn) and
\texttt{bridge\_search\_CIRCULAR\_retracted.py}.

\section*{Open questions}

The uniform two-defect floor reduces to a single arithmetic input:
\textbf{Hypothesis~$\beta$ with $\beta>1/2$ on the interior family
$U,V\ge1$} (\S\ref{sec:beta}). It is open, and it is an effective question.
Everything else proved here is a per-configuration certificate.

Two further items are recorded as limitations rather than results. First, the
prior-art check of \S\ref{sec:lit} is not exhaustive: it is a targeted audit of
the works nearest to this setting, not a systematic survey, and absence from it
is not evidence of novelty. Second, no citation here rests on an abstract
alone---every work cited has been read in full, except three classical items
cited through Lagarias's survey rather than their originals, as marked.

\section*{Provenance}

The results were produced twice, by two independently written implementations
of the same specification: one via the affine identity of
\S\ref{sec:class}--\S\ref{sec:r1}, one via a backward recursion on the word
that derives the modulus $2^{S+1}$ with no input from the author. They agree on
every shared quantity, including the closed form (0 mismatches out of
$19\,773$) and all floor values. A third, separately produced package was
obtained and executed, reproducing its reported correlation $r=-0.179$, merge
supremum $+0.3209$, balanced medians $0.337/0.671/1.015$ and $\beta^{*}=1/2$.

Three claims made during development were withdrawn and are recorded here
because the corrected forms are used above. (i) A bridge-style search reported
positive margins by estimating its transfer constant from the measured
$m(w)$; that is circular, since the transfer condition rearranges to the
conclusion it was meant to establish (\S\ref{sec:neg}). (ii) The two
representatives of \S\ref{sec:tworun} were claimed to differ in height by at
most one bit; they differ by one \emph{modulus}, but by up to $9.43$ bits of
height on the grids tested. (iii) The gain from Theorem~\ref{thm:fc}\,(c) was
claimed to be $O(1)$; Proposition~\ref{prop:unbounded} shows it is unbounded.

\begin{thebibliography}{9}
\bibitem{akhtar2026}
M.~A.~K. Akhtar, \emph{Compression Spectra of Collatz Orbits: Sharp Residual
Ordering, Exact Valuation Cylinders, and Residue-Threshold Sparsity},
preprint, 19 August 2026, \texttt{doi:10.5281/zenodo.22015248};
SSRN~7314419. \emph{Full text consulted.}

\bibitem{bohmsontacchi}
C.~B\"ohm and G.~Sontacchi, Atti Accad. Naz. Lincei Rend. Cl. Sci. Fis. Mat.
Natur. \textbf{64} (1978), 260--264. \emph{Verified first-hand.}

\bibitem{dejesus-eoc}
E. De~Jes\'us, \emph{A Global Occupation Conjecture for the Accelerated $3x+1$
Map: Residue--Time Duality, Divergent-Orbit Sparsity, and the Pointwise
Realizer Problem}, Revision~8, 4 October 2026,
\texttt{doi:10.5281/zenodo.23146031}. \emph{Verified at source.}

\bibitem{dejesus-erratum}
E. De~Jes\'us, \emph{The Least-Realizer Minimum for Confined Words: An Erratum
to the Terras--Everett Residue Class\dots}, 11 August 2026,
\texttt{doi:10.5281/zenodo.21889704}. \emph{Verified at source.}

\bibitem{everett1977}
C.~J. Everett, \emph{Iteration of the number theoretic function
$f(2n)=n$, $f(2n+1)=3n+2$}, Advances in Math. \textbf{25} (1977), 42--45.
\emph{Cited via \cite{lagarias2010}; original not consulted.}

\bibitem{hh1997}
L.~Halbeisen and N.~Hungerb\"uhler, \emph{Optimal bounds for the length of
rational Collatz cycles}, Acta Arith. \textbf{78} (1997), no.~3, 227--239.
\emph{Verified first-hand.}

\bibitem{kayadibi2026}
S.~Y. Kayadibi, \emph{A Modular Classification of Pre-Descent Resistance in
Accelerated Odd Collatz Dynamics}, 18 June 2026,
\texttt{doi:10.5281/zenodo.20745800}. \emph{Full text consulted.}

\bibitem{lagarias1985}
J.~C. Lagarias, \emph{The $3x+1$ problem and its generalizations}, Amer. Math.
Monthly \textbf{92} (1985), 3--23. \emph{Theorem~B cited via
\cite{lagarias2010}; original not consulted.}

\bibitem{lagarias2010}
J.~C. Lagarias, \emph{The $3x+1$ problem: an overview}, in: The Ultimate
Challenge: The $3x+1$ Problem, Amer. Math. Soc., 2010, pp.~3--29;
arXiv:2111.02635. \emph{Full text consulted.}

\bibitem{lagarias1990}
J.~C. Lagarias, \emph{The set of rational cycles for the $3x+1$ problem},
Acta Arith. \textbf{56} (1990), 33--53. \emph{Not re-verified at source.}

\bibitem{rozierterracol}
O.~Rozier and C.~Terracol, \emph{Paradoxical behavior in Collatz sequences},
arXiv:2502.00948v5 (17 May 2026); Discrete Math. \textbf{349} (2026).
\emph{Full text consulted.}

\bibitem{terras1976}
R.~Terras, \emph{A stopping time problem on the positive integers}, Acta Arith.
\textbf{30} (1976), 241--252. \emph{Cited via \cite{lagarias2010};
original not consulted.}
\end{thebibliography}

\section*{Acknowledgement of AI assistance}

AI assistance (Claude) was used for verification and drafting. Conceptual
direction and final responsibility are the author's.

\section*{Licensing}

This paper is released under CC BY 4.0. The accompanying code is released under
the MIT License.

\end{document}
